\section{Why Evaluation-Awareness Probes Are Confounded}
The existing diagnostic compares a pooled probe trained on all four cells with one-format probes trained on benchmark or casual cells. The pooled probe has both more rows and direct exposure to both target formats. Thus its advantage is not a clean test of format-conditioned purpose coding. We retain this diagnostic as development evidence only and make equal-N training the first inferential gate.

Probe selectivity and control-task work motivates matched controls \citep{hewitt2019control,belinkov2021probing}. Causal probing work also distinguishes decodability from behavioral use \citep{elazar2020amnesic}; our intervention claims follow that distinction.

Recent work makes the adjacent methodological space clear: format-confound controls and evaluation-context transfer motivate the present factorial design \citep{sahoo2026format,devbunova2026evaluation}, while activation transport and representation interventions provide causal tools \citep{rodriguez2025act,wu2024reft,geiger2024das}. A closely related structure$\times$register factorial study \citep{santana2026relational} reinforces that interaction terms can matter, but does not estimate our purpose$\times$format residual-space $A/B/C$ decomposition or compare held-out interventions with natural matched downstream targets. Our novelty claim is therefore the combination, not factorial interactions or activation intervention in isolation.
